% ==============================================================================
% SECTION 4 — RESULTS
% ==============================================================================
\section{Results}
\label{sec:results}

\subsection{Strategy Validation: Cooperation Propensity}

Table~\ref{tab:cooperation} reports the normalised propensity to cooperate
for the Default prompt without noise, analogous to Table~3 in Willis et al.

\begin{table}[h]
\caption{Normalised cooperation propensity (Default prompt, no noise).
         Rows: row-player attitude; columns: opponent attitude.}
\label{tab:cooperation}
\small
\setlength{\tabcolsep}{4pt}
\begin{tabular}{@{}llccc@{}}
\toprule
\textbf{Model} & \textbf{Att.} & \textbf{vs A} & \textbf{vs C} & \textbf{vs N} \\
\midrule
\multirow{3}{*}{Claude 4.6}
  & A & 0.000 & 0.297 & 0.286 \\
  & C & 0.298 & 0.991 & 0.992 \\
  & N & 0.287 & 0.993 & 1.000 \\
\midrule
\multirow{3}{*}{Gemini 2.5 Flash}
  & A & 0.082 & 0.087 & 0.087 \\
  & C & 0.088 & 1.000 & 1.000 \\
  & N & 0.088 & 1.000 & 1.000 \\
\midrule
\multirow{3}{*}{Gemini 3.1 Pro}
  & A & 0.413 & 0.629 & 0.629 \\
  & C & 0.630 & 0.999 & 0.999 \\
  & N & 0.630 & 0.999 & 0.999 \\
\midrule
\multirow{3}{*}{GPT-5.4 Mini}
  & A & 0.006 & 0.025 & 0.023 \\
  & C & 0.026 & 1.000 & 0.940 \\
  & N & 0.024 & 0.940 & 0.883 \\
\bottomrule
\end{tabular}
\end{table}

All four models exhibit the expected attitude separation: aggressive strategies
cooperate markedly less than cooperative or neutral strategies.
Gemini~3.1~Pro is the notable exception, with aggressive strategies cooperating
at 41\% even against other aggressive strategies---substantially higher than any
model tested by Willis et al. By contrast, Gemini~2.5~Flash and GPT-5.4~Mini
generate highly committed aggressive strategies (cooperation rates near zero
against all opponents), while Claude~4.6 closely mirrors Claude~3.5~Sonnet's
pattern: aggressive strategies cooperate at $\approx 0$\% against other
aggressors but rise to $\approx 29$\% against cooperative opponents.

\subsection{Head-to-Head Payoffs and Differential Capabilities}

Table~\ref{tab:payoffs} presents the normalised mean payoffs across all
25 matches for each attitude pairing (Default prompt, no noise), together with
the \ICD (Eq.~\ref{eq:icd}).

\begin{table*}[t]
\caption{Normalised head-to-head payoffs (Default prompt, no noise) and
         Index of Differential Capabilities (\ICD). Lower \ICD indicates
         a larger cooperative advantage.}
\label{tab:payoffs}
\small
\begin{tabular}{@{}llccccccc@{}}
\toprule
\textbf{Model} & \textbf{Prompt} &
  \multicolumn{3}{c}{\textbf{Aggressive payoff vs.}} &
  \multicolumn{3}{c}{\textbf{Cooperative payoff vs.}} &
  \textbf{\ICD} \\
\cmidrule(lr){3-5}\cmidrule(lr){6-8}
 &  & A & C & N & A & C & N & \\
\midrule
\multirow{3}{*}{Claude 4.6}
  & Default & 1.000 & 1.893 & 1.861 & 1.890 & 2.982 & 2.988 & 0.605 \\
  & Prose   & 1.242 & 1.917 & 1.784 & 1.507 & 3.000 & 2.836 & 0.673 \\
  & Refine  & 2.432 & 2.820 & 2.683 & 2.772 & 2.997 & 2.924 & 0.913 \\
\midrule
\multirow{3}{*}{Gemini 2.5 Flash}
  & Default & 1.232 & 1.244 & 1.244 & 1.239 & 3.000 & 3.000 & 0.514 \\
  & Prose   & 1.290 & 1.796 & 1.551 & 0.985 & 2.574 & 2.086 & 0.822 \\
  & Refine  & 1.486 & 2.288 & 2.022 & 1.470 & 2.998 & 2.635 & 0.816 \\
\midrule
\multirow{3}{*}{Gemini 3.1 Pro}
  & Default & 1.992 & 2.400 & 2.400 & 2.397 & 2.998 & 2.998 & 0.809 \\
  & Prose   & 2.009 & 1.971 & 1.972 & 1.702 & 2.979 & 2.858 & 0.789 \\
  & Refine  & 2.159 & 2.405 & 2.433 & 2.259 & 2.685 & 2.617 & 0.925 \\
\midrule
\multirow{3}{*}{GPT-5.4 Mini}
  & Default & 1.017 & 1.079 & 1.073 & 1.074 & 3.000 & 2.900 & 0.454 \\
  & Prose   & 1.982 & 2.395 & 2.423 & 2.394 & 2.914 & 2.937 & 0.825 \\
  & Refine  & 1.066 & 1.427 & 1.252 & 1.343 & 2.846 & 2.305 & 0.577 \\
\bottomrule
\end{tabular}
\end{table*}

\ICD values range from 0.454 (GPT-5.4~Mini Default) to 0.925 (Gemini~3.1~Pro~Refine)
and vary substantially across prompts. Cooperative and neutral attitudes achieve
near-mutual-cooperation payoffs ($\approx 3.0$) in most model--prompt combinations,
consistent with the original findings. GPT-5.4~Mini Default's \ICD of 0.454
means aggressive strategies achieve less than half the payoff of cooperative ones;
at the other extreme, Gemini~3.1~Pro~Refine (0.925) and Claude~4.6~Refine (0.913)
approach payoff parity between aggressive and cooperative agents. The Self-Refine
prompt increases \ICD in all four models, replicating the original finding that
Self-Refine narrows the aggressive--cooperative gap. For Claude~4.6~Refine,
this effect is particularly pronounced: aggressive strategies achieve payoffs
of 2.43--2.82, approaching cooperative strategies at 2.77--3.00.

Compared to Willis et al.'s Claude~3.5~Sonnet (Table~4 therein), Claude~4.6
shows systematically higher payoffs for aggressive strategies under Refine,
suggesting improved aggressive strategy generation in the newer version.
GPT-5.4~Mini Default, however, shows weaker aggressive performance than
ChatGPT-4o~Default (1.05 vs.\ 1.81 average), indicating that model scaling
does not uniformly improve aggressive capability.

\subsection{Evolutionary Equilibria}

Table~\ref{tab:moran} reports the main Moran process results: the proportion
of 500 runs converging to each equilibrium attitude, for all four population
conditions. The Willis et al.\ reference values are shown for comparison.

\begin{table*}[t]
\caption{Moran process equilibrium proportions (\%A / \%C / \%N) across
         four population conditions. Prior probability indicates the
         baseline for each initial ratio.
         Reference values from Willis et al.\ \cite{Willis2025_llm_ipd}
         (Table~6) are shown in the bottom rows.
         Bold: values highlighted in the text discussion.}
\label{tab:moran}
\footnotesize
\setlength{\tabcolsep}{3pt}
\begin{tabular}{@{}llcccc@{}}
\toprule
\textbf{Model} & \textbf{Prompt} &
  \textbf{4:4:4 clean} &
  \textbf{4:4:4 noise} &
  \textbf{8:2:2 clean} &
  \textbf{8:2:2 noise} \\
\cmidrule(lr){3-3}\cmidrule(lr){4-4}\cmidrule(lr){5-5}\cmidrule(lr){6-6}
 & & (prior: 33/33/33) & (prior: 33/33/33) & (prior: 67/17/17) & (prior: 67/17/17) \\
\midrule
\multirow{3}{*}{Claude 4.6}
  & Default & 2/49/\textbf{49} & 24/40/36 & 12/47/41 & 60/23/17 \\
  & Prose   & 5/\textbf{48}/47 & 28/37/36 & 32/38/30 & 64/17/19 \\
  & Refine  & 23/\textbf{39}/38 & 29/41/30 & 59/22/20 & 62/23/15 \\
\midrule
\multirow{3}{*}{Gemini 2.5 Flash}
  & Default & 2/\textbf{49}/48 & 20/39/41 & 24/38/38 & 61/23/16 \\
  & Prose   & 24/\textbf{47}/30 & 40/28/32 & 65/19/16 & 77/9/14 \\
  & Refine  & 21/\textbf{47}/32 & 36/32/32 & 56/22/22 & 73/13/13 \\
\midrule
\multirow{3}{*}{Gemini 3.1 Pro}
  & Default & 14/42/\textbf{45} & 31/35/34 & 36/32/32 & 63/20/17 \\
  & Prose   & 19/\textbf{43}/38 & 36/29/34 & 60/21/20 & 70/12/17 \\
  & Refine  & 27/\textbf{40}/33 & 32/37/31 & 62/18/20 & 66/18/16 \\
\midrule
\multirow{3}{*}{GPT-5.4 Mini}
  & Default & 2/\textbf{53}/45 & 14/47/39 & 22/38/40 & 52/25/23 \\
  & Prose   & 12/\textbf{45}/43 & 31/34/35 & 46/28/26 & 66/16/18 \\
  & Refine  & 4/\textbf{70}/26 & 25/42/33 & 28/46/26 & 60/19/20 \\
\midrule
\midrule
\multirow{3}{*}{\textit{ChatGPT-4o}$^\dagger$}
  & Default & 14/53/33 & 16/42/42 & 66/19/17 & 59/20/21 \\
  & Prose   & 13/38/49 & 23/41/36 & 35/27/38 & 60/18/22 \\
  & Refine  & 19/48/33 & 28/38/34 & 49/30/21 & 63/19/18 \\
\multirow{3}{*}{\textit{Claude 3.5 Sonnet}$^\dagger$}
  & Default & 4/49/47  & 15/37/48 & 36/24/40 & 41/20/39 \\
  & Prose   & 14/42/44 & 17/33/50 & 41/30/29 & 61/26/13 \\
  & Refine  & 16/51/33 & 37/34/29 & 50/22/28 & 60/18/22 \\
\bottomrule
\end{tabular}
\vspace{2pt}\\
\small$^\dagger$From Willis et al.\ \cite{Willis2025_llm_ipd}, Table~6.
\end{table*}

\paragraph{Balanced, noiseless condition (4:4:4 clean).}
This condition provides the most direct comparison with the original benchmark.
Ten of twelve model--prompt combinations favour cooperative equilibria as
the plurality outcome ($\PC > \PA$ and $\PC > \PN$), confirming a widespread
cooperative bias. The strongest cooperative equilibria are achieved
by GPT-5.4~Mini~Refine (70\%), GPT-5.4~Mini~Default (53\%), and
Claude~4.6~Prose (48\%).
The two exceptions are both cases where the \emph{Neutral} attitude achieves the
plurality: Gemini~3.1~Pro~Default (14\%A/42\%C/\textbf{45\%N}) and
Claude~4.6~Default (2\%A/49\%C/\textbf{49\%N}), the latter a virtual
cooperative--neutral tie. Notably, neither exception is driven by
aggressive dominance; in both cases aggressive equilibria remain well below
the 33\% prior. The Neutral plurality for Gemini~3.1~Pro likely reflects that its
aggressive strategies cooperate at unusually high rates (Table~\ref{tab:cooperation},
41\% vs.\ own attitude), blurring the boundary between aggressive and neutral
behaviour and allowing neutral strategies to benefit from mutual cooperation without
defection penalties.

Comparing within-lineage successors:
Claude~4.6~Default achieves 2\%A/49\%C/49\%N, virtually identical to
Claude~3.5~Sonnet~Default at 4\%A/49\%C/47\%N. This is the closest
generational replication in the dataset.
GPT-5.4~Mini~Default achieves 2\%A/53\%C/45\%N, far lower aggressive than
ChatGPT-4o~Default (14\%A) and now favouring a cooperative plurality, suggesting
improved cooperative bias in the OpenAI lineage for this condition.

\paragraph{Balanced, noisy condition (4:4:4 noise).}
Noise consistently shifts equilibria toward aggression across all twelve
model--prompt combinations; all twelve pairs show
$\PA^{\text{noise}} > \PA^{\text{clean}}$.
Cooperative equilibria degrade under noise in every combination but one:
Claude~4.6~Refine is marginally higher under noise ($\Dnoise = -1$pp;
see Section~\ref{sec:noise} and Table~\ref{tab:noise}), the only case
in which cooperative equilibria do not decline.
Claude~4.6 is the most noise-robust on average ($\Dnoise = 6$pp), while
GPT-5.4~Mini~Refine shows the largest single degradation:
cooperative equilibria fall from 70\% to 42\% ($\Dnoise = 28$pp).

\paragraph{Biased, noiseless condition (8:2:2 clean).}
With aggressive agents as the initial majority, the prior for aggressive
equilibria is 67\%. Below this threshold indicates resistance to aggressive
takeover; above it indicates susceptibility. All twelve model--prompt
combinations finish below the prior, meaning cooperative and neutral
strategies resist takeover even when outnumbered. Resistance is strongest
under Default prompts---Claude~4.6 (12\%A), GPT-5.4~Mini (22\%A) and
Gemini~2.5~Flash (24\%A) all finish far below the 67\% prior. The closest
to the prior is Gemini~2.5~Flash~Prose at 65\%A, still short of the initial
aggressive proportion.

\paragraph{Biased, noisy condition (8:2:2 noise).}
In the most adverse condition (aggressive majority with noise), all twelve
model--prompt combinations converge to aggressive equilibria at rates of
52--77\%. Contrary to a naive reading of the aggressive-majority prior (67\%),
only 3 of 12 combinations exceed this threshold---Gemini~2.5~Flash~Prose (77\%),
Gemini~2.5~Flash~Refine (73\%) and Gemini~3.1~Pro~Prose (70\%). The remaining
9 of 12 fall at or below the prior, indicating that the Moran process does not
amplify aggressive dominance beyond the initial seeding for most model--prompt
combinations, even in this worst-case scenario. The most resistant are
GPT-5.4~Mini~Default (52\%) and Claude~4.6~Default (60\%), which finish below
the prior.

\subsection{Noise Sensitivity ($\Dnoise$)}
\label{sec:noise}

Table~\ref{tab:noise} summarises $\Dnoise$ (Eq.~\ref{eq:dnoise})
for the balanced condition, providing a direct measure of noise robustness.

\begin{table}[h]
\caption{Noise sensitivity $\Dnoise$ (balanced population, 4:4:4).
         Positive values indicate degradation of cooperative equilibria under noise.}
\label{tab:noise}
\small
\begin{tabular}{@{}lllll@{}}
\toprule
\textbf{Model} & \textbf{Default} & \textbf{Prose} & \textbf{Refine} & \textbf{Avg.} \\
\midrule
Claude 4.6        & 9  & 12 & $-1$ & 6  \\
Gemini 2.5 Flash  & 10 & 19 & 14   & 15 \\
Gemini 3.1 Pro    & 6  & 14 & 3    & 8  \\
GPT-5.4 Mini      & 7  & 11 & 28   & 15 \\
\midrule
\textit{Claude 3.5 Sonnet}$^\dagger$ & 12 & 9 & 17 & 13 \\
\textit{ChatGPT-4o}$^\dagger$         & 11 & $-3$ & 10 & 6 \\
\bottomrule
\end{tabular}
\vspace{2pt}\\
\small$^\dagger$Derived from Willis et al.\ \cite{Willis2025_llm_ipd}, Table~6.
\end{table}

Claude~4.6 stands out as the most noise-robust on average ($\Dnoise = 6$pp),
followed by Gemini~3.1~Pro (8pp). Gemini~2.5~Flash and GPT-5.4~Mini are the
most noise-sensitive overall (avg.\ 15pp each). Claude~4.6's average (6pp)
is markedly lower than Claude~3.5~Sonnet's (13pp), directionally consistent
with improved noise robustness in the newer model---though, as the
cross-study test in the Discussion shows, this gap is not statistically
robust once the predecessor's smaller sample size is propagated.
